\def\handoutsay{false}\def\handoutfill{false}
\documentclass[12pt]{article}
\newcommand{\handoutrunning}{Lecture 11: Linear Models I}
\providecommand{\handoutsay}{true}
\usepackage[T1]{fontenc}
\usepackage{lmodern,microtype}
\usepackage[a4paper,margin=15mm,headheight=14pt,headsep=5mm]{geometry}
\usepackage{amsmath,amssymb,amsthm,mathtools}
\usepackage{enumitem,booktabs,needspace}
\usepackage{pifont}
\usepackage[dvipsnames]{xcolor}
\usepackage{fancyhdr}
\usepackage{ifthen}
\usepackage[hidelinks]{hyperref}
\hypersetup{pdfauthor={Chang-Ye Tu}}
\pagestyle{fancy}
\fancyhf{}
\fancyhead[L]{\small Linear Algebra}
\fancyhead[R]{\small\handoutrunning}
\fancyfoot[C]{\small\thepage}
\renewcommand{\headrulewidth}{0.3pt}
\setlength{\parindent}{0pt}
\setlength{\parskip}{4pt plus 1pt minus 1pt}
\setlength{\emergencystretch}{1.5em}
\usepackage{etoolbox}
\newcommand{\handoutdisplayskips}{\setlength{\abovedisplayskip}{7pt plus 2pt minus 2pt}%
  \setlength{\belowdisplayskip}{7pt plus 2pt minus 2pt}}
\apptocmd{\normalsize}{\handoutdisplayskips}{}{\PackageWarning{handout}{could not patch normalsize}}
\handoutdisplayskips
\setlist[enumerate]{label=(\alph*),leftmargin=*,itemsep=2pt,topsep=3pt}
\setlist[itemize]{leftmargin=*,itemsep=2pt,topsep=3pt}
\newcommand{\R}{\mathbb R}
\newcommand{\C}{\mathbb C}
\newcommand{\iu}{\mathrm{i}}
\newcommand{\contradiction}{\mathord{\text{\large\ding{56}}}}
\newcommand{\one}{\mathbf 1}
\DeclareMathOperator{\tr}{tr}
\DeclareMathOperator{\col}{col}
\DeclareMathOperator{\diag}{diag}
\DeclareMathOperator{\rank}{rank}
\DeclareMathOperator\prb{\mathsf{P}}
\DeclareMathOperator\expc{\mathsf{E}}
\DeclareMathOperator\var{var}
\DeclareMathOperator\cov{cov}
\DeclareMathOperator\cor{cor}
\DeclareMathOperator\argmin{argmin}
\DeclareMathOperator\argmax{argmax}
\DeclareMathOperator{\Span}{span}
\theoremstyle{definition}
\newtheorem{theorem}{Theorem}
\newtheorem{proposition}[theorem]{Proposition}
\newtheorem{lemma}[theorem]{Lemma}
\newtheorem{corollary}[theorem]{Corollary}
\newtheorem{definition}[theorem]{Definition}
\newtheorem{example}[theorem]{Example}
\newtheorem{remark}[theorem]{Remark}
\newtheorem{fact}{Fact}
\renewcommand{\thefact}{F\arabic{fact}}
\AddToHook{cmd/subsection/before}{\par\Needspace{9\baselineskip}}
\AddToHook{env/theorem/before}{\par\Needspace{4\baselineskip}}
\AddToHook{env/proposition/before}{\par\Needspace{4\baselineskip}}
\AddToHook{env/lemma/before}{\par\Needspace{4\baselineskip}}
\AddToHook{env/definition/before}{\par\Needspace{4\baselineskip}}
\AddToHook{env/example/before}{\par\Needspace{4\baselineskip}}
\definecolor{saybg}{gray}{0.86}
\NewDocumentEnvironment{say}{+b}
  {\ifthenelse{\equal{\handoutsay}{true}}{\par\Needspace{2\baselineskip}\noindent\colorbox{saybg}{\parbox{\dimexpr\linewidth-2\fboxsep\relax}%
   {\small\sffamily\ding{43}\enspace\textbf{Say.}\enspace#1}}\par}{}}
  {}
\providecommand{\handoutfill}{false}
\newsavebox{\boardbox}
\newlength{\boardht}
\definecolor{boardrule}{gray}{0.55}
\makeatletter
\newcommand{\boardstart}{\par\if@inlabel\leavevmode\par\fi}
\makeatother
\NewDocumentEnvironment{board}{o +b}
  {\boardstart\IfValueT{#1}{\noindent#1\par\nopagebreak}%
   \noindent\begin{lrbox}{\boardbox}\begin{minipage}{\linewidth}#2\end{minipage}\end{lrbox}%
   \setlength{\boardht}{\dimexpr\ht\boardbox+\dp\boardbox\relax}%
   \ifthenelse{\equal{\handoutfill}{true}}%
     {\setlength{\boardht}{1.4\boardht}%
      {\setlength{\fboxsep}{0pt}\setlength{\fboxrule}{0.4pt}\color{boardrule}%
       \framebox[\dimexpr\linewidth-2\fboxrule\relax]{\rule{0pt}{\boardht}}}}%
     {\ifthenelse{\equal{\handoutsay}{true}}%
        {\llap{\color{boardrule}\rule[-\dp\boardbox]{1.5pt}{\boardht}\hspace{4pt}}}{}%
      \usebox{\boardbox}}%
   \par}
  {}
\newcommand{\blank}[1]{\ifthenelse{\equal{\handoutfill}{true}}{\underline{\phantom{#1}}}{#1}}
\newcommand{\lec}[1]{Lecture~#1}
\newcommand{\handouttitle}[2]{%
  \begin{center}
  {\LARGE\sffamily\bfseries #2\par}\vspace{3pt}
  {\sffamily Linear Algebra \textperiodcentered\ \lec{#1}\par}
  \ifthenelse{\equal{\handoutsay}{true}}{\vspace{3pt}{\small\sffamily\colorbox{saybg}{\ding{43}\enspace\textbf{Say.}}\enspace
    teacher's copy: the shaded boxes are not in the students' copy.\par}}{}
  \ifthenelse{\equal{\handoutfill}{true}}{\vspace{3pt}{\small\sffamily Fill-in copy: complete the framed spaces in
    class.\par}}{}
  \end{center}}

\hypersetup{pdftitle={Linear Algebra, Lecture 11: Linear Models I}}
\newcommand{\boardtop}{\vspace*{-\abovedisplayskip}}
\DeclareMathOperator{\row}{row}
\begin{document}
\handouttitle{11}{Linear Models I}

\textbf{Last time}: \(z\sim N_n(\mu,\Sigma)\Rightarrow Az+b\sim N_m(A\mu+b,A\Sigma A')\) (\lec{10}, Theorem~3);
\(z\sim N_n(\mu,\sigma^2I_n)\), \(A'=A=A^2\), \(\rank A=r\geqslant1\), \(A\mu=0\) \(\Rightarrow\)
\(z'Az/\sigma^2\sim\chi^2_r\) (\lec{10}, Theorem~15). Earlier results are cited where they are used. \(e\) is the vector of
residuals; unit vectors are written as columns of \(I\).

\subsection*{1. The linear model and least squares}

\begin{definition}\label{h-def-model}
In the \textbf{linear model} \(y=X\beta+\varepsilon\), the \textbf{design matrix} \(X\in\R^{n\times p}\) is a known
constant matrix, \(\beta\in\R^p\) is an unknown constant vector of \textbf{parameters}, the \textbf{errors}
\(\varepsilon\in\R^n\) form a random vector, and the \textbf{response} \(y\in\R^n\) is observed; \(x_i'\) is row \(i\) of
\(X\). In \textbf{multiple regression}, \(p=k+1\), \(\beta=(\beta_0,\dots,\beta_k)'\) and
\(x_i'=(1,x_{i1},\dots,x_{ik})\), with \(x_{ij}\) the value of the \(j\)th explanatory variable on observation \(i\): the
first column \(\one_n=(1,\dots,1)'\) of \(X\) (\lec{1}, Definition~14) is the column of the \textbf{intercept}
\(\beta_0\). The assumptions:
(A1) \(\expc(y)=X\beta\); (A2) \(\expc(\varepsilon)=0\); (A3) \(\var(\varepsilon)=\sigma^2I_n\) (\(\sigma^2>0\)); (A4)
\(\rank X=p\). Results hold for \(X\) of any rank unless (A4) is assumed. A \textbf{least-squares estimate} of \(\beta\) is a
\(\hat\beta\in\R^p\) that minimizes \(S(\beta)=\lVert y-X\beta\rVert^2=\sum_i(y_i-x_i'\beta)^2\); the \(p\) equations
\(X'X\beta=X'y\) are the \textbf{normal equations}.
\end{definition}
Here \(\varepsilon\), \(y\) are random vectors (\lec{9}, Definition~1: vectors of random variables). By \lec{9},
Theorem~3(b) (\(\expc(Az+b)=A\expc(z)+b\), \(\var(Az+b)=A\var(z)A'\); \(A=I_n\), \(b=X\beta\)):
\(\expc(y)=X\beta+\expc(\varepsilon)\), \(\var(y)=\var(\varepsilon)\). So (A1) \(\iff\) (A2), and (A3) gives
\(\var(y)=\sigma^2I_n\). (A4) needs \(n\geqslant p\) (\lec{4}, Theorem~12(a): \(\rank X\leqslant\min\{n,p\}\)). An
\textbf{estimator} \(T=g(y)\) of a parameter \(\theta\) is a random variable, \(g\) involving no unknown parameter; it is
\textbf{unbiased} if \(\expc(T)=\theta\) for all values of the parameters (\lec{9}, Definition~7).

\Needspace{14\baselineskip}
\begin{theorem}\label{h-thm-normal}
Let \(X\in\R^{n\times p}\) and \(y\in\R^n\).
\begin{enumerate}
\item The normal equations \(X'X\beta=X'y\) have a solution.
\item \(X'X\) is invertible \(\iff\) \(\rank X=p\). Then \(\hat\beta=(X'X)^{-1}X'y\) is the only solution.
\item If \(\hat\beta\) is a solution, the solutions are the vectors \(\hat\beta+z\), \(z\in\mathcal N(X)\); all give the same
\(X\hat\beta\). If \(\rank X<p\), there are infinitely many solutions.
\end{enumerate}
\end{theorem}
\begin{proof}
\begin{board}[(a) \(\mathcal C(A)=\{Ax\}\), \(\mathcal N(A)=\{x:Ax=0\}\) (\lec{4}, Definition~7):]
\boardtop
\begin{align*}
X'y&\in\mathcal C(X')=\mathcal C(X'X)&&\text{(\lec{4}, Theorem~16(c))}\\
&\implies X'X\beta=X'y\text{ has a solution}&&\text{(\lec{4}, Corollary~13(a))}
\end{align*}
\end{board}
\begin{board}[(b) \(X'X\in\R^{p\times p}\):]
\boardtop
\begin{align*}
X'X\text{ invertible}&\iff\rank(X'X)=p\iff\rank X=p\qquad\text{(\lec{4}, Theorems~12(d), 16(b))}\\
X'X\beta=X'y&\implies\beta=(X'X)^{-1}X'X\beta=(X'X)^{-1}X'y,\text{ and this }\beta\text{ is a solution}
\end{align*}
\end{board}
\begin{board}[(c) If \(Ax_p=b\), the solutions of \(Ax=b\) are the \(x_p+z\), \(Az=0\) (\lec{2}, Theorem~22(a)):]
\boardtop
\begin{align*}
\text{solutions }\hat\beta+z,\ z&\in\mathcal N(X'X)=\mathcal N(X)\qquad\text{(\lec{4}, Theorem~16(a))}\\
X(\hat\beta+z)&=X\hat\beta+Xz=X\hat\beta\\
\rank X<p&\implies\dim\mathcal N(X)=p-\rank X\geqslant1\qquad\text{(\lec{4}, Theorem~11)}\\
&\implies\text{some }z\ne0\text{ in }\mathcal N(X)\qquad\text{(\lec{4}, Theorem~4: \(\dim\{0\}=0\))}\\
\hat\beta+t_1z&=\hat\beta+t_2z\implies t_1=t_2:\ \text{the solutions }\hat\beta+tz\ (t\in\R)\text{ are distinct}
\end{align*}
\end{board}
\nopagebreak
\end{proof}

\Needspace{16\baselineskip}
\begin{theorem}\label{h-thm-ls}
Let \(X\in\R^{n\times p}\) and \(y\in\R^n\).
\begin{enumerate}
\item \(\hat\beta\) is a least-squares estimate \(\iff\) \(X'X\hat\beta=X'y\). For every solution \(\hat\beta\) of the normal
equations: \(S(\beta)=S(\hat\beta)+\lVert X(\hat\beta-\beta)\rVert^2\) (\(\beta\in\R^p\)).
\item All least-squares estimates give the same vector of \textbf{fitted values} \(\hat y=X\hat\beta\), the orthogonal
projection of \(y\) onto \(\mathcal C(X)\); the vector of \textbf{residuals} \(e=y-\hat y\) lies in \(\mathcal C(X)^\perp\).
\item If \(\rank X=p\): \(\hat\beta=(X'X)^{-1}X'y\) is the only least-squares estimate, \(\hat y=Hy\) and \(e=My\), with the
\textbf{hat matrix} \(H=X(X'X)^{-1}X'\) (\lec{8}, Theorem~3: \(P=A(A'A)^{-1}A'\), here \(A=X\)) and the
\textbf{residual maker} \(M=I_n-H\).
\end{enumerate}
\end{theorem}
\begin{proof}
\(y\perp\mathcal C(A)\iff A'y=0\) (\lec{6}, Theorem~5(a)); the orthogonal projection of \(y\) onto \(W\): the
\(\hat w\in W\) with \(y-\hat w\perp W\) (\lec{8}, Definition~1); \(\lVert x\rVert=0\iff x=0\) (\lec{6}, Definition~1).
\begin{board}[(a) \(\Leftarrow\): \(X'X\hat\beta=X'y\); \(X\beta\in\mathcal C(X)\) for \(\beta\in\R^p\).
\(\Rightarrow\): \(\beta^*\) a least-squares estimate, \(\hat\beta\) a solution (Theorem~\ref{h-thm-normal}(a)):]
\boardtop
\begin{align*}
X'(y-X\hat\beta)&=0\implies y-X\hat\beta\perp\mathcal C(X)\qquad\text{(\lec{6}, Theorem~5(a))}\\
&\implies X\hat\beta\text{ is the projection of }y\text{ onto }\mathcal C(X)\qquad\text{(\lec{8}, Definition~1)}\\
\lVert y-X\beta\rVert^2&=\lVert y-X\hat\beta\rVert^2+\lVert X\hat\beta-X\beta\rVert^2\implies S(\beta)\geqslant S(\hat\beta)
\qquad\text{(\lec{8}, Theorem~2)}\\
\lVert X(\hat\beta-\beta^*)\rVert^2&=S(\beta^*)-S(\hat\beta)\leqslant0\implies X\beta^*=X\hat\beta\qquad\text{(\(\beta^*\)
minimizes \(S\); \lec{6}, Definition~1)}\\
X'X\beta^*&=X'X\hat\beta=X'y
\end{align*}
\end{board}
(b) By the last two lines of (a), every least-squares estimate \(\beta^*\) gives \(X\beta^*=X\hat\beta\), the projection
of the first two lines; and \(e=y-X\hat\beta\perp\mathcal C(X)\), so \(e\in\mathcal C(X)^\perp\) (\lec{6}, Definition~4:
\(W^\perp=\{y:y\perp W\}\)).

(c) By Theorem~\ref{h-thm-normal}(b) and (a), \((X'X)^{-1}X'y\) is the only least-squares estimate; then
\(\hat y=X(X'X)^{-1}X'y=Hy\) and \(e=y-Hy=My\).
\end{proof}

\begin{proposition}\label{h-prop-residuals}
Let \(\hat\beta\) be a least-squares estimate, \(\hat y=X\hat\beta\), \(e=y-\hat y\).
(a) \(X'e=0\). (b) If \(\one_n\in\mathcal C(X)\): \(\one'e=0\), so \(\hat y\) and \(y\) have the same mean \(\bar y\)
(the mean of \(x\in\R^n\) is \(\bar x=\frac1n\one'x\), \lec{1}, Example~16).
(c) \(\hat y'e=0\).
\end{proposition}
\begin{proof}
(a) \(X'(y-X\hat\beta)=0\) is the normal equations (Theorem~\ref{h-thm-ls}(a)).

(b) \(\one=Xc\): \(\one'e=c'X'e=0\) by (a), and \(\one'\hat y=\one'y-\one'e=\one'y\).

(c) \(\hat y'e=\hat\beta'X'e=0\) by (a).
\end{proof}

\Needspace{16\baselineskip}
\begin{example}\label{h-exm-ls}
Let \(X=[\one_4\ (2,0,2,0)'\ (3,3,1,1)']\in\R^{4\times3}\) (\(k=2\)) and \(y=(9,3,5,3)'\): the \(A\) and \(y\) of
\lec{8}, Example~6, where (with \(H=P\))
\[
X'X=\begin{pmatrix}4&4&8\\4&8&8\\8&8&20\end{pmatrix},\quad
(X'X)^{-1}=\frac14\begin{pmatrix}6&-1&-2\\-1&1&0\\-2&0&1\end{pmatrix},\quad
(X'X)^{-1}X'=\frac14\begin{pmatrix}-2&0&2&4\\1&-1&1&-1\\1&1&-1&-1\end{pmatrix}
\]
and \(Hy=(8,4,6,2)'\), \((I-H)y=(1,-1,-1,1)'\). Find \(\hat\beta\), \(\hat y\), \(e\), \(S(\hat\beta)\) and
\(S((1,2,2)')\).
\end{example}
\textit{Solution.}
\begin{board}[\(X'X\) is invertible, so \(\rank X=3=p\) (Theorem~\ref{h-thm-normal}(b)): (A4) holds; \(X'y\) has the
entries \(\col_j(X)'y\) (\lec{1}, Theorem~4(d), Definition~7):]
\boardtop
\begin{align*}
X'y&=\begin{pmatrix}9+3+5+3\\18+10\\27+9+5+3\end{pmatrix}=\begin{pmatrix}20\\28\\44\end{pmatrix},\qquad
\hat\beta=(X'X)^{-1}X'y=\frac14\begin{pmatrix}-18+0+10+12\\9-3+5-3\\9+3-5-3\end{pmatrix}=\begin{pmatrix}1\\2\\1\end{pmatrix}\\
X'X\hat\beta&=(4+8+8,\ 4+16+8,\ 8+16+20)'=(20,28,44)'=X'y
\end{align*}
\end{board}
\begin{board}[Theorem~\ref{h-thm-ls}; Proposition~\ref{h-prop-residuals} for the checks; the identity of
Theorem~\ref{h-thm-ls}(a) for \(S((1,2,2)')\):]
\boardtop
\begin{align*}
\hat y&=X\hat\beta=(1+4+3,\ 1+0+3,\ 1+4+1,\ 1+0+1)'=(8,4,6,2)'=Hy\\
e&=y-\hat y=(1,-1,-1,1)'\\
S(\hat\beta)&=e'e=1+1+1+1=4>0:\ X\beta=y\text{ has no solution}\\
X'e&=(1-1-1+1,\ 2-2,\ 3-3-1+1)'=0,\qquad\one'e=0,\qquad\hat y'e=8-4-6+2=0\\
S((1,2,2)')&=4+\lVert X(0,0,-1)'\rVert^2=4+(9+9+1+1)=24
\end{align*}
\end{board}

\subsection*{2. Rank-deficient designs and generalized inverses}

\Needspace{18\baselineskip}
\begin{theorem}\label{h-thm-oneway}
Let \(n_1,\dots,n_g\geqslant1\), \(n=n_1+\dots+n_g\); \(n\) observations in \(g\) groups (first the \(n_1\) of group
\(1\), then the \(n_2\) of group \(2\), and so on); the \textbf{indicator} \(d_l\in\R^n\) of group \(l\) has entry \(1\) in
the rows of group \(l\) and \(0\) elsewhere. The model \(y_i=\mu+\alpha_l+\varepsilon_i\) (observation \(i\) in group
\(l\)) is \(y=X\beta+\varepsilon\), \(\beta=(\mu,\alpha_1,\dots,\alpha_g)'\), \(X=[\one_n\ d_1\ \cdots\ d_g]\in\R^{n\times(g+1)}\).
Let \(X_1=[d_1\ \cdots\ d_g]\), \(X_2=[\one_n\ d_2\ \cdots\ d_g]\in\R^{n\times g}\).
\begin{enumerate}
\item \(X_1\) has full column rank, and \(X=X_1S\) with \(S=[\one_g\ I_g]\in\R^{g\times(g+1)}\) of full row rank.
\item \(\rank X=g\) and \(\mathcal N(X)=\Span\{(-1,1,\dots,1)'\}\).
\item \(X_1=XT_1\), \(X_2=XT_2\), with \(T_1,T_2\in\R^{(g+1)\times g}\) the matrix \(I_{g+1}\) without its first, resp.\
second, column. \(X_1\), \(X_2\) have full column rank \(g\), and \(\mathcal C(X_1)=\mathcal C(X_2)=\mathcal C(X)\).
\end{enumerate}
\end{theorem}
\begin{proof}
\(A\in\R^{m\times n}\) has full column rank \(\iff\) its columns are independent \(\iff\) \(\mathcal N(A)=\{0\}\), and full
row rank \(\iff\) \(\mathcal C(A)=\R^m\) (\lec{4}, Theorem~12(b),(c)).
\begin{board}[(a) No observation lies in two groups, each lies in one:]
\boardtop
\begin{align*}
d_l'd_m&=0\ (l\ne m),\qquad d_l'd_l=n_l\\
c_1d_1+\dots+c_gd_g&=0\implies c_mn_m=0\ (\text{multiply by }d_m')\implies c_m=0\\
&\implies X_1\text{ has full column rank}&&\text{(\lec{4}, Theorem~12(b))}\\
d_1+\dots+d_g&=\one_n\implies X_1\one_g=\one_n&&\text{(\lec{1}, Theorem~4(b))}\\
&\implies X_1S=[X_1\one_g\ \ X_1I_g]=X&&\text{(\lec{1}, Theorem~4(a))}\\
\mathcal C(S)&\supseteq\mathcal C(I_g)=\R^g:\ S\text{ full row rank}&&\text{(\lec{4}, Theorem~12(c))}
\end{align*}
\end{board}
\begin{board}[(b) \(X_1\) has full column rank ((a)), so \(\rank(X_1B)=\rank B\), \(\mathcal N(X_1B)=\mathcal N(B)\)
(\lec{4}, Theorem~15(b)):]
\boardtop
\begin{align*}
\rank X&=\rank(X_1S)=\rank S=g\qquad\text{(\lec{4}, Theorem~15(b); (a))}\\
\mathcal N(X)&=\mathcal N(X_1S)=\mathcal N(S)\qquad\text{(\lec{4}, Theorem~15(b))}\\
S(c_0,c_1,\dots,c_g)'&=c_0\one_g+(c_1,\dots,c_g)'=0\iff c_l=-c_0\ (l=1,\dots,g)\\
&\iff(c_0,c_1,\dots,c_g)'=-c_0(-1,1,\dots,1)'
\end{align*}
\end{board}
(c)
\begin{itemize}[itemsep=0pt,topsep=1pt]
\item \(X\col_j(T_1)\) is a column of \(X\) other than the first (\lec{1}, Theorem~4(a),(b)), so \(XT_1=X_1\); likewise
\(XT_2=X_2\).
\item \(\one_n=d_1+\dots+d_g\) and \(d_1=\one_n-d_2-\dots-d_g\): each column of \(X_2\) lies in \(\mathcal C(X_1)\) and each
column of \(X_1\) in \(\mathcal C(X_2)\); so \(\mathcal C(X_2)=\mathcal C(X_1)\), of dimension \(g\) by (a), and \(X_2\),
with \(g\) columns, has full column rank (\lec{4}, Definition~10: \(\rank A=\dim\mathcal C(A)\)).
\item \(\mathcal C(X_1)=\mathcal C(XT_1)\subseteq\mathcal C(X)\) (\lec{4}, Theorem~15(a)), both of dimension \(g\) ((b)):
equal (\lec{4}, Theorem~5(a): a subspace \(U\subseteq W\) with \(\dim U=\dim W\) is \(W\)).
\end{itemize}
\end{proof}
So the intercept together with all \(g\) indicators (the \textbf{dummy variable trap}) makes \(X'X\) singular
(Theorem~\ref{h-thm-normal}(b)); dropping the intercept or one indicator (\(d_1\) in Theorem~\ref{h-thm-oneway}(c); any \(d_l\) after renumbering the
groups) gives a design of full column rank with the
same column space, hence the same fitted values (Theorem~\ref{h-thm-ls}(b)).

\Needspace{20\baselineskip}
\begin{example}\label{h-exm-oneway}
Six plants of three types weigh \(5,7,6\) (type \(1\)), \(3,5\) (type \(2\)) and \(2\) (type \(3\)): \(y=(5,7,6,3,5,2)'\),
the model of Theorem~\ref{h-thm-oneway} with \(g=3\), \((n_1,n_2,n_3)=(3,2,1)\), and (entries
\(\col_j(X)'\col_l(X)\), \(\col_j(X)'y\): \lec{1}, Theorem~4(d), Definition~7)
\[
X=\begin{pmatrix}1&1&0&0\\1&1&0&0\\1&1&0&0\\1&0&1&0\\1&0&1&0\\1&0&0&1\end{pmatrix},\qquad
X'X=\begin{pmatrix}6&3&2&1\\3&3&0&0\\2&0&2&0\\1&0&0&1\end{pmatrix},\qquad
X'y=\begin{pmatrix}28\\18\\8\\2\end{pmatrix}
\]
Find all least-squares estimates and \(\hat y\), and the estimates with \(X_1\) and with \(X_2\).
\end{example}
\textit{Solution.}
\begin{board}[\(\rank X=3\), \(\mathcal N(X)=\Span\{(-1,1,1,1)'\}\) (Theorem~\ref{h-thm-oneway}(b)): \(X'X\) is singular
(Theorem~\ref{h-thm-normal}(b)); all solutions, that is, all least-squares estimates, by Theorems~\ref{h-thm-normal}(c)
and~\ref{h-thm-ls}(a):]
\boardtop
\begin{align*}
X'X(0,6,4,2)'&=(18+8+2,\ 18,\ 8,\ 2)'=(28,18,8,2)'=X'y\\
\hat\beta&=(0,6,4,2)'+t(-1,1,1,1)'=(-t,\ 6+t,\ 4+t,\ 2+t)'\quad(t\in\R)\\
X\hat\beta&=(6,6,6,4,4,2)'\ \text{for all }t:\ \text{the group means }18/3,\ 8/2,\ 2/1\\
\hat\alpha_1-\hat\alpha_2&=(6+t)-(4+t)=2\ \text{for all }t;\qquad\hat\alpha_1=6+t\text{ changes with }t
\end{align*}
\end{board}
\begin{board}[\(X_1\), \(X_2\) have full column rank and \(\mathcal C(X_1)=\mathcal C(X_2)=\mathcal C(X)\)
(Theorem~\ref{h-thm-oneway}(c)): one solution each (Theorem~\ref{h-thm-normal}(b)), the same \(\hat y\)
(Theorem~\ref{h-thm-ls}(b)):]
\boardtop
\begin{align*}
X_1'X_1&=\diag(3,2,1),\qquad X_1'y=(18,8,2)',\qquad\hat\beta=(X_1'X_1)^{-1}X_1'y=(6,4,2)'\\
X_2'X_2&=\begin{pmatrix}6&2&1\\2&2&0\\1&0&1\end{pmatrix},\qquad X_2'y=\begin{pmatrix}28\\8\\2\end{pmatrix},\qquad
X_2'X_2\begin{pmatrix}6\\-2\\-4\end{pmatrix}=\begin{pmatrix}36-4-4\\12-4\\6-4\end{pmatrix}=X_2'y\\
\hat y&=6d_1+4d_2+2d_3=6\,\one_6-2d_2-4d_3=(6,6,6,4,4,2)'
\end{align*}
\end{board}
\((6,4,2)'\) and \((6,-2,-4)'\) are the coordinates of \(\hat y\) in the bases \(d_1,d_2,d_3\) and \(\one_6,d_2,d_3\) of
\(\mathcal C(X)\) (\lec{3}, Theorem~15: the coordinates in a basis are unique). In this \textbf{baseline coding}, the coefficient of \(\one_6\) is the value of group
\(1\), and those of \(d_2\), \(d_3\) are the differences of groups \(2\), \(3\) from it. \(T_1(6,4,2)'=(0,6,4,2)'\) and
\(T_2(6,-2,-4)'=(6,0,-2,-4)'\) are the solutions with \(t=0\) and \(t=-6\).

\begin{definition}\label{h-def-ginverse}
Let \(A\in\R^{m\times n}\). A matrix \(G\in\R^{n\times m}\) with \(AGA=A\) is a \textbf{generalized inverse}
(\textbf{g-inverse}) of \(A\); \(A^-\) denotes any g-inverse of \(A\). If \(y=Az\in\mathcal C(A)\), then
\(A(Gy)=AGAz=Az=y\): \(x=Gy\) solves \(Ax=y\). If \(A\in\R^{n\times n}\) is invertible, \(A^{-1}\) is its only
g-inverse: \(AGA=A\) gives \(G=A^{-1}(AGA)A^{-1}=A^{-1}AA^{-1}=A^{-1}\).
\end{definition}

\Needspace{16\baselineskip}
\begin{theorem}\label{h-thm-submatrix}
Let \(A\in\R^{m\times n}\), \(\rank A=r\geqslant1\). For rows \(i_1<\dots<i_r\) and columns \(j_1<\dots<j_r\), let the
submatrix \(A_*\in\R^{r\times r}\) (\((s,t)\) entry \(a_{i_sj_t}\)) be invertible, and let \(G\in\R^{n\times m}\) have
\((j_s,i_t)\) entry equal to the \((s,t)\) entry of \(A_*^{-1}\) (\(s,t=1,\dots,r\)) and all other entries \(0\). Then
\(G\) is a g-inverse of \(A\). In particular, if
\(A=\left(\begin{smallmatrix}A_{11}&A_{12}\\A_{21}&A_{22}\end{smallmatrix}\right)\) with \(A_{11}\in\R^{r\times r}\)
invertible, then \(\left(\begin{smallmatrix}A_{11}^{-1}&0\\0&0\end{smallmatrix}\right)\in\R^{n\times m}\) is a g-inverse
of \(A\).
\end{theorem}
\begin{proof}
\begin{board}[Selection matrices \(U=[\col_{i_1}(I_m)\ \cdots\ \col_{i_r}(I_m)]\) and
\(V=[\col_{j_1}(I_n)\ \cdots\ \col_{j_r}(I_n)]\); the \((s,t)\) entry of \(U'AV\) is
\(\col_{i_s}(I_m)'A\col_{j_t}(I_n)=a_{i_sj_t}\) (\lec{1}, Theorem~4(d), Proposition~11(a)):]
\boardtop
\begin{align*}
U'AV&=A_*,\qquad B=AV\in\R^{m\times r},\qquad\mathcal C(B)\subseteq\mathcal C(A)&&\text{(\lec{4}, Theorem~15(a))}\\
Bz&=0\implies A_*z=U'Bz=0\implies z=0\\
&\implies B\text{ has independent columns}&&\text{(\lec{4}, Theorem~12(b))}\\
\dim\mathcal C(A)&=r\implies\text{these }r\text{ columns form a basis of }\mathcal C(A)&&\text{(\lec{4}, Theorem~5(d))}\\
&\implies A=BT,\ T\in\R^{r\times n}&&\text{(\lec{1}, Theorem~4(a),(b))}\\
U'A&=U'BT=A_*T\implies T=A_*^{-1}U'A\\
A(VA_*^{-1}U')A&=BA_*^{-1}U'A=BT=A\\
VA_*^{-1}U'&=\textstyle\sum_{s,t}(A_*^{-1})_{st}\col_{j_s}(I_n)\col_{i_t}(I_m)'=G&&\text{(\lec{1}, Theorem~4(d))}
\end{align*}
\end{board}
\nopagebreak
\end{proof}
The case \(A_*=A_{11}\) is \(i_s=j_s=s\) (\(s=1,\dots,r\)). So: pick an invertible \(r\times r\) submatrix, invert it, put
\(A_*^{-1}\) (not transposed) into rows \(j_1,\dots,j_r\) and columns \(i_1,\dots,i_r\) of the \(n\times m\) zero
matrix.

\Needspace{11\baselineskip}
\begin{theorem}\label{h-thm-hat}
Let \(X\in\R^{n\times p}\), \(r=\rank X\), \(G\) a g-inverse of \(X'X\), and \(H=XGX'=X(X'X)^-X'\), \(M=I_n-H\) (the
\textbf{hat matrix} and the \textbf{residual maker}).
\begin{enumerate}
\item \(XGX'X=X\) and \(X'XGX'=X'\); \(\beta^0=GX'y\) is a solution of the normal equations.
\item For every \(y\in\R^n\), \(Hy\) is the orthogonal projection of \(y\) onto \(\mathcal C(X)\). So \(H\) does not depend
on the choice of \(G\), and \(H=X(X'X)^{-1}X'\) if \(r=p\).
\item \(H'=H=H^2\), \(HX=X\), \(X'H=X'\); \(M'=M=M^2\), \(MX=0\), \(HM=MH=0\).
\item \(\rank H=\tr H=r\) and \(\rank M=\tr M=n-r\).
\item For every least-squares estimate \(\hat\beta\): \(\hat y=Hy\), \(e=My\), and the \textbf{error sum of squares} is
\(\mathrm{SSE}=S(\hat\beta)=e'e=y'My=y'y-\hat\beta'X'y\).
\end{enumerate}
\end{theorem}
\begin{proof}
\begin{board}[(a) \(X'X(GX'X-I_p)=X'XGX'X-X'X=0\) (Definition~\ref{h-def-ginverse}), so the columns of \(GX'X-I_p\)
lie in \(\mathcal N(X'X)\) (\lec{1}, Theorem~4(a)):]
\boardtop
\begin{align*}
\mathcal N(X'X)&=\mathcal N(X)\implies X(GX'X-I_p)=0:\ XGX'X=X&&\text{(\lec{4}, Theorem~16(a))}\\
X'XG'X'X&=(X'XGX'X)'=(X'X)'=X'X&&\text{(\lec{1}, Theorem~8)}\\
&\implies G'\text{ is a g-inverse of }X'X\\
XG'X'X&=X\ (\text{line 1 for }G')\implies X'XGX'=(XG'X'X)'=X'&&\text{(\lec{1}, Theorem~8)}
\end{align*}
\(X'y\in\mathcal C(X'X)\) (Theorem~\ref{h-thm-normal}(a)), so \(X'X(GX'y)=X'y\) (Definition~\ref{h-def-ginverse}).
\end{board}
\begin{board}[(b) By (a), \lec{6}, Theorem~5(a) and \lec{8}, Definition~1 (the projection is unique):]
\boardtop
\begin{align*}
Hy&=X(GX'y)\in\mathcal C(X),\qquad X'(y-Hy)=X'y-X'XGX'y=0\\
&\implies Hy\text{ is the projection of }y\text{ onto }\mathcal C(X)\\
G_1,G_2\text{ g-inverses}&\implies XG_1X'y=XG_2X'y\ (\text{all }y)\\
&\implies XG_1X'=XG_2X'\qquad\text{(\lec{1}, Proposition~11(b))}
\end{align*}
If \(r=p\): \(X'X\) is invertible (Theorem~\ref{h-thm-normal}(b)), \((X'X)^{-1}\) is its only g-inverse
(Definition~\ref{h-def-ginverse}), and \(H=X(X'X)^{-1}X'\).
\end{board}
\begin{board}[(c) \(G'\) is a g-inverse of \(X'X\) ((a)), so \(H=XG'X'\) ((b)):]
\boardtop
\begin{align*}
H'&=(XGX')'=XG'X'=H&&\text{(\lec{1}, Theorem~8)}\\
HX&=XGX'X=X,\qquad X'H=X'H'=(HX)'=X'&&\text{((a); \lec{1}, Theorem~8)}\\
H^2&=(HX)GX'=XGX'=H\\
M'&=I-H'=M,\qquad M^2=I-2H+H^2=M\\
MX&=X-HX=0,\qquad HM=H-H^2=0=MH
\end{align*}
\end{board}
\begin{board}[(d) \(H\), \(M\) are symmetric and idempotent ((c)):]
\boardtop
\begin{align*}
\mathcal C(H)&=\mathcal C(X(GX'))\subseteq\mathcal C(X)=\mathcal C(HX)\subseteq\mathcal C(H)&&\text{(\lec{4}, Theorem~15(a))}\\
\rank H&=\dim\mathcal C(X)=r&&\text{(\lec{4}, Definition~10)}\\
\rank H&=\tr H,\qquad\rank M=\tr M&&\text{(\lec{8}, Theorem~11)}\\
\tr M&=\tr I_n-\tr H=n-r&&\text{(\lec{1}, Theorem~10(a))}
\end{align*}
\end{board}
(e) \(\hat y=X\hat\beta\) is the projection of \(y\) onto \(\mathcal C(X)\) (Theorem~\ref{h-thm-ls}(b)), which is \(Hy\) by
(b); so \(e=y-Hy=My\). By (c), \(S(\hat\beta)=e'e=y'M'My=y'My\); and \(y'Hy=y'X\hat\beta=\hat\beta'X'y\)
(\lec{1}, Theorem~8), so \(y'My=y'y-\hat\beta'X'y\).
\end{proof}

\Needspace{12\baselineskip}
\begin{example}\label{h-exm-ginverse}
In Example~\ref{h-exm-oneway}, find a g-inverse \(G\) of \(X'X\) by Theorem~\ref{h-thm-submatrix}, \(\beta^0=GX'y\) and
SSE.
\end{example}
\textit{Solution.}
\begin{board}[\(\rank(X'X)=\rank X=3\) (\lec{4}, Theorem~16(b)); rows and columns \(2,3,4\) of \(X'X\) give
\(A_*=\diag(3,2,1)\), \(A_*^{-1}=\diag(\frac13,\frac12,1)\) (\lec{1}, Proposition~13(b); \lec{2}, Theorem~15):]
\boardtop
\begin{align*}
G&=\diag(0,\tfrac13,\tfrac12,1)&&\text{(Theorem~\ref{h-thm-submatrix})}\\
\beta^0&=GX'y=(0,\ 18/3,\ 8/2,\ 2)'=(0,6,4,2)'&&\text{(Theorem~\ref{h-thm-hat}(a))}\\
\mathrm{SSE}&=y'y-\beta^{0\prime}X'y=148-(0+108+32+4)=4&&\text{(Theorem~\ref{h-thm-hat}(e))}\\
e&=y-X\beta^0=(-1,1,0,-1,1,0)',\qquad e'e=4
\end{align*}
Here \(y'y=25+49+36+9+25+4=148\), and \(\beta^0\) is the solution with \(t=0\) in Example~\ref{h-exm-oneway}.
\end{board}

\Needspace{13\baselineskip}
\subsection*{3. The Gauss--Markov theorem}

\Needspace{11\baselineskip}
\begin{theorem}\label{h-thm-unbiased}
Assume (A4), and let \(\hat\beta=(X'X)^{-1}X'y\).
\begin{enumerate}
\item \(\hat\beta=\beta+(X'X)^{-1}X'\varepsilon\); under (A2), \(\expc(\hat\beta)=\beta\).
\item Under (A3), \(\var(\hat\beta)=\sigma^2(X'X)^{-1}\).
\end{enumerate}
\end{theorem}
\begin{proof}
\begin{board}[(a) \(y=X\beta+\varepsilon\):]
\boardtop
\begin{align*}
\hat\beta&=(X'X)^{-1}X'(X\beta+\varepsilon)=\beta+(X'X)^{-1}X'\varepsilon&&\text{(\lec{1}, Theorem~3(a))}\\
\expc(\hat\beta)&=\beta+(X'X)^{-1}X'\expc(\varepsilon)=\beta&&\text{(\lec{9}, Theorem~3(b), \(z=\varepsilon\), \(b=\beta\); (A2))}
\end{align*}
\end{board}
\begin{board}[(b) \(((X'X)^{-1})'=(X'X)^{-1}\) (\lec{1}, Theorem~8; \lec{2}, Corollary~12); by (a):]
\[
\var(\hat\beta)=(X'X)^{-1}X'(\sigma^2I_n)X(X'X)^{-1}=\sigma^2(X'X)^{-1}\qquad\text{(\lec{9}, Theorem~3(b); (A3))}
\]
\end{board}
\nopagebreak
\end{proof}

\Needspace{16\baselineskip}
\begin{theorem}[Gauss--Markov]\label{h-thm-gm}
Assume (A1)--(A4), and let \(\hat\beta=(X'X)^{-1}X'y\). Let \(L\in\R^{p\times n}\) be constant, and \(\tilde\beta=Ly\) a
\textbf{linear unbiased estimator} of \(\beta\): \(\expc(Ly)=\beta\) for every \(\beta\in\R^p\). Then
\(\var(\tilde\beta)-\var(\hat\beta)\) is PSD (\lec{8}, Definition~13: symmetric, \(x'Ax\geqslant0\) for all \(x\)), and it
is \(0\) only if \(\tilde\beta=\hat\beta\). So
\(\var(a'\tilde\beta)\geqslant\var(a'\hat\beta)\) for every \(a\in\R^p\); in particular
\(\var(\tilde\beta_j)\geqslant\var(\hat\beta_j)\) for each entry. \(\hat\beta\) is the \textbf{BLUE} (best linear unbiased
estimator) of \(\beta\).
\end{theorem}
\begin{proof}
\begin{board}[\(LX=I_p\) (\lec{9}, Theorem~3(b), (A1)); the difference \(D\):]
\boardtop
\begin{align*}
\expc(Ly)&=L\expc(y)=LX\beta=\beta\ \text{for all }\beta\\
\beta&=\col_j(I_p)\ (j=1,\dots,p)\implies LX=I_p&&\text{(\lec{1}, Theorem~4(b))}\\
D&=L-(X'X)^{-1}X'\implies DX=LX-I_p=0&&\text{(\lec{1}, Theorem~3(a))}\\
X'D'&=(DX)'=0&&\text{(\lec{1}, Theorem~8)}
\end{align*}
\end{board}
\begin{board}[The variances (\lec{9}, Theorem~3(b), (A3)); \(DX=0\), \(X'D'=0\), \((X'X)^{-1}\) symmetric
(\lec{2}, Corollary~12); \((D')'D'\) is PSD (\lec{8}, Theorem~15(b)):]
\boardtop
\begin{align*}
\var(\tilde\beta)&=L(\sigma^2I_n)L'=\sigma^2LL'\\
LL'&=(X'X)^{-1}X'X(X'X)^{-1}+(X'X)^{-1}X'D'\\
&\qquad+DX(X'X)^{-1}+DD'=(X'X)^{-1}+DD'&&\text{(\lec{1}, Theorems~3(a), 8)}\\
\var(\tilde\beta)-\var(\hat\beta)&=\sigma^2DD'=\sigma^2(D')'D'\text{ is PSD}&&\text{(Theorem~\ref{h-thm-unbiased}(b))}\\
\sigma^2DD'=0&\implies\tr\bigl((D')'D'\bigr)=0\implies D'=0\implies\tilde\beta=\hat\beta&&\text{(\lec{1}, Theorem~10(c))}\\
\var(a'\tilde\beta)-\var(a'\hat\beta)&=\sigma^2a'DD'a=\sigma^2\lVert D'a\rVert^2\geqslant0&&\text{(\lec{6}, Definition~1)}
\end{align*}
\(a=\col_j(I_p)\) gives \(a'\tilde\beta=\tilde\beta_j\): \(\var(\tilde\beta_j)\geqslant\var(\hat\beta_j)\).
\end{board}
\nopagebreak
\end{proof}

\Needspace{12\baselineskip}
\begin{example}\label{h-exm-blue}
In Example~\ref{h-exm-ls}, under (A1)--(A4): show that \(c'y=\frac12(y_1-y_2)\) has \(\expc(c'y)=\beta_1\) for every
\(\beta\), and compare its variance with that of the entry \(\hat\beta_1\) of the BLUE \(\hat\beta\).
\end{example}
\textit{Solution.}
\begin{board}[\(x_1'=(1,2,3)\), \(x_2'=(1,0,3)\), \(c=\frac12(1,-1,0,0)'\), \(a=(0,1,0)'\); (A1), (A3), \lec{9},
Theorem~3(b); Theorem~\ref{h-thm-unbiased}(b) for \(\var(\hat\beta_1)\); \(M'=M\), \(MX=0\)
(Theorem~\ref{h-thm-hat}(c)):]
\boardtop
\begin{align*}
\expc(y_1-y_2)&=(x_1-x_2)'\beta=(0,2,0)\beta=2\beta_1,\qquad\expc(c'y)=\beta_1=a'\beta\\
\var(c'y)&=\sigma^2c'c=\sigma^2/2;\qquad\hat\beta_1=\tfrac14(y_1-y_2+y_3-y_4),\quad\var(\hat\beta_1)=\sigma^2\cdot\tfrac14\\
X'c&=a,\qquad Hc=X(X'X)^{-1}X'c=X(X'X)^{-1}a=\tfrac14(1,-1,1,-1)'\\
u&=c-Hc=Mc=\tfrac14(1,-1,-1,1)'\\
c'&=a'(X'X)^{-1}X'+u',\qquad u'X=c'MX=0\\
\var(c'y)-\var(\hat\beta_1)&=\frac{\sigma^2}2-\frac{\sigma^2}4=\sigma^2\lVert u\rVert^2=\sigma^2\cdot\tfrac4{16}
\end{align*}
\(\hat\beta_1\) is row \(2\) of \((X'X)^{-1}X'y\). For \(y=(9,3,5,3)'\) the two estimates are \(\frac12(9-3)=3\) and
\(\hat\beta_1=2\).
\end{board}

\Needspace{14\baselineskip}
\begin{proposition}\label{b13.21}
Let \(X\in\R^{n\times p}\), \(y\in\R^n\), and \(W\in\R^{n\times n}\) PD (\lec{8}, Definition~13: symmetric,
\(x'Wx>0\) for \(x\ne0\)). The minimizers of \(Q(\beta)=(y-X\beta)'W(y-X\beta)\) (\(\beta\in\R^p\)) are the solutions of
\(X'WX\beta=X'Wy\). If \(\rank X=p\), then \(X'WX\) is PD, and the unique minimizer of \(Q\) is the \textbf{weighted
least-squares estimate} \(\hat\beta_W=(X'WX)^{-1}X'Wy\). For \(W=\diag(w_1,\dots,w_n)\) with all \(w_i>0\),
\(Q(\beta)=\sum_iw_i(y_i-x_i'\beta)^2\).
\end{proposition}
\begin{proof}
\begin{board}[\(T=W^{1/2}\) (\lec{10}, Proposition~12: \(W^{1/2}\), \(W^{-1/2}\) symmetric, \(W^{1/2}W^{1/2}=W\),
\(W^{1/2}W^{-1/2}=I\)), so \(T'=T\), \(T'T=W\), and \(T\) is invertible (\lec{2}, Theorem~15):]
\boardtop
\begin{align*}
Q(\beta)&=(y-X\beta)'T'T(y-X\beta)=\lVert Ty-TX\beta\rVert^2&&\text{(\lec{1}, Theorem~8; \lec{6}, Definition~1)}\\
&\text{minimizers: }(TX)'TX\beta=(TX)'Ty&&\text{(Theorem~\ref{h-thm-ls}(a) for \(TX\), \(Ty\))}\\
&\iff X'WX\beta=X'Wy&&(T'T=W)
\end{align*}
\end{board}
\begin{board}[\(\rank X=p\):]
\boardtop
\begin{align*}
\rank(TX)&=\rank X=p&&\text{(\lec{4}, Theorem~15(d): \(\rank(UAV)=\rank A\), \(U,V\) invertible)}\\
X'WX&=(TX)'(TX)\text{ is PD}&&\text{(\lec{8}, Theorem~15(b): \(A'A\) is PD \(\iff\) \(\rank A=p\))}
\end{align*}
The only minimizer is \(\hat\beta_W=\bigl((TX)'TX\bigr)^{-1}(TX)'Ty=(X'WX)^{-1}X'Wy\) (Theorem~\ref{h-thm-ls}(c) for
\(TX\), \(Ty\)).
\end{board}
For \(W=\diag(w_1,\dots,w_n)\), all \(w_i>0\): \(W'=W\), \(x'Wx=\sum_iw_ix_i^2>0\) for \(x\ne0\), so \(W\) is PD, and
\(Q(\beta)=\sum_iw_iu_i^2\) for \(u=y-X\beta\) (\lec{1}, Definition~2, Proposition~13(a)).
\end{proof}

\Needspace{16\baselineskip}
\subsection*{4. Estimating \(\sigma^2\); sums of squares}

\Needspace{12\baselineskip}
\begin{theorem}\label{h-thm-sigma}
Assume (A2), (A3), and \(\rank X=r<n\); let \(G\) be a g-inverse of \(X'X\) and \(M=I_n-XGX'\) (Theorem~\ref{h-thm-hat}).
Then \(\expc(\mathrm{SSE})=(n-r)\sigma^2\), so
\(\hat\sigma^2=\mathrm{SSE}/(n-r)\) is an unbiased estimator of \(\sigma^2\). Under (A4), \(r=p\), and \(\hat\sigma^2\) is
the \textbf{mean squared error} \(s^2=\mathrm{MSE}=\mathrm{SSE}/(n-p)\) (\(=\mathrm{SSE}/(n-k-1)\) in multiple regression).
\end{theorem}
\begin{proof}
\begin{board}[\(y=X\beta+\varepsilon\); \lec{9}, Theorem~6: \(\expc(z'Az)=\tr(A\Sigma)+\mu'A\mu\) for \(\expc(z)=\mu\),
\(\var(z)=\Sigma\); here \(z=\varepsilon\), \(\mu=0\) (A2), \(\Sigma=\sigma^2I_n\) (A3):]
\boardtop
\begin{align*}
e&=My=MX\beta+M\varepsilon=M\varepsilon&&\text{(Theorem~\ref{h-thm-hat}(c),(e))}\\
\mathrm{SSE}&=e'e=\varepsilon'M'M\varepsilon=\varepsilon'M\varepsilon&&\text{(Theorem~\ref{h-thm-hat}(c))}\\
\expc(\mathrm{SSE})&=\tr(M\sigma^2I_n)+0'M0=\sigma^2\tr M=\sigma^2(n-r)
&&\text{(\lec{1}, Theorem~10(a); Theorem~\ref{h-thm-hat}(d))}\\
\expc(\hat\sigma^2)&=\tfrac1{n-r}\expc(\mathrm{SSE})=\sigma^2&&\text{(\(n-r>0\); \lec{9}, Theorem~3(b))}
\end{align*}
\end{board}
\nopagebreak
\end{proof}
Example~\ref{h-exm-ls}: \(\mathrm{SSE}=S(\hat\beta)=4\), \(n-r=4-3=1\), so \(s^2=4\), and
\(s^2(X'X)^{-1}=\left(\begin{smallmatrix}6&-1&-2\\-1&1&0\\-2&0&1\end{smallmatrix}\right)\) estimates \(\var(\hat\beta)\)
(Theorem~\ref{h-thm-unbiased}(b)). Example~\ref{h-exm-ginverse}: \(\hat\sigma^2=4/(6-3)=\frac43\).

\begin{definition}\label{h-def-ss}
For a least-squares fit, with \(\hat y\) and \(e=y-\hat y\): the \textbf{total}, \textbf{regression} and \textbf{error sums
of squares} are \(\mathrm{SST}=\sum_i(y_i-\bar y)^2\), \(\mathrm{SSR}=\sum_i(\hat y_i-\bar y)^2\),
\(\mathrm{SSE}=\sum_i(y_i-\hat y_i)^2=e'e\). If \(\one_n\in\mathcal C(X)\) and \(\mathrm{SST}>0\), the
\textbf{coefficient of determination} is \(R^2=\mathrm{SSR}/\mathrm{SST}\).
\end{definition}

\Needspace{13\baselineskip}
\begin{theorem}\label{h-thm-anova}
If \(\one_n\in\mathcal C(X)\): \(\mathrm{SST}=\mathrm{SSR}+\mathrm{SSE}\), with \(\mathrm{SST}=\lVert C_ny\rVert^2=y'C_ny\)
and \(\mathrm{SSR}=\lVert C_n\hat y\rVert^2\) (\(C_n=I_n-\frac1n\one\one'\), \lec{1}, Definition~14); if
\(\mathrm{SST}>0\), \(R^2=1-\mathrm{SSE}/\mathrm{SST}\in[0,1]\); and \(y'y=n\bar y^2+\mathrm{SSR}+\mathrm{SSE}\).
\end{theorem}
\begin{proof}
\(C_n'=C_n=C_n^2\), \(C_n\one=0\) (\lec{1}, Proposition~15); \(C_nx=x-\bar x\one\) and
\(\sum_i(x_i-\bar x)^2=x'C_nx\) for \(x\in\R^n\) (\lec{1}, Example~16).
\begin{board}[\(\hat y\) and \(y\) have the mean \(\bar y\), \(\one'e=0\) (Proposition~\ref{h-prop-residuals}(b)):]
\boardtop
\begin{align*}
\mathrm{SST}&=\lVert C_ny\rVert^2=y'C_ny,\qquad\mathrm{SSR}=\lVert C_n\hat y\rVert^2&&\text{(\lec{1}, Proposition~15, Example~16)}\\
C_ny&=C_n\hat y+C_ne=C_n\hat y+e&&(C_ne=e-\bar e\one=e)\\
(C_n\hat y)'e&=\hat y'C_ne=\hat y'e=0&&\text{(Proposition~\ref{h-prop-residuals}(c))}\\
\mathrm{SST}&=\lVert C_n\hat y\rVert^2+\lVert e\rVert^2=\mathrm{SSR}+\mathrm{SSE}&&\text{(\lec{6}, Theorem~2: Pythagoras)}\\
\mathrm{SST}>0&\implies R^2=1-\mathrm{SSE}/\mathrm{SST},\quad0\leqslant\mathrm{SSR}\leqslant\mathrm{SST}&&(\mathrm{SSE}\geqslant0)\\
y&=\bar y\one+C_ny,\qquad(\bar y\one)'C_ny=\bar y(C_n\one)'y=0&&\text{(\lec{1}, Proposition~15)}\\
y'y&=n\bar y^2+\lVert C_ny\rVert^2=n\bar y^2+\mathrm{SSR}+\mathrm{SSE}&&\text{(\lec{6}, Theorem~2)}
\end{align*}
\end{board}
\nopagebreak
\end{proof}

\Needspace{9\baselineskip}
\begin{example}\label{h-exm-anova}
In Example~\ref{h-exm-ls} (\(\one_4\) is a column of \(X\)), find SST, SSR, SSE, \(R^2\) and the partition of \(y'y\).
\end{example}
\textit{Solution.}
\begin{board}[\(\bar y=\frac14(9+3+5+3)=5\), \(\hat y=(8,4,6,2)'\) (Example~\ref{h-exm-ls}); Theorem~\ref{h-thm-anova}:]
\boardtop
\begin{align*}
C_4y&=(4,-2,0,-2)',\qquad C_4\hat y=(3,-1,1,-3)',\qquad e=(1,-1,-1,1)',\qquad C_4y=C_4\hat y+e\\
\mathrm{SST}&=16+4+0+4=24,\quad\mathrm{SSR}=9+1+1+9=20,\quad\mathrm{SSE}=4:\quad24=20+4,\quad R^2=\tfrac{20}{24}=\tfrac56\\
y'y&=81+9+25+9=124=4\cdot25+20+4=100+20+4
\end{align*}
\end{board}

\bigskip
\textbf{Practice.} The Exercises and Problems of the textbook, Chapters~12 and~13: generalized inverses, least squares
and the normal equations, the hat matrix, rank-deficient designs, the Gauss--Markov theorem, the estimation of
\(\sigma^2\), and sums of squares.

\end{document}
